% problem_id: S001-theorem-coherent-algebraic
% source_id: S001 (Stacks Project)
% source_file: quot.tex
% statement_locator: quot.tex lines 1448--1455
% proof_locator: quot.tex lines 1457--1529
% env: theorem
% id_note: label 在本来源内唯一
% status: complete (statement + author proof verbatim + reason + locators)
%
\begin{theorem}[Algebraicity of the stack of coherent sheaves; flat case]
\label{theorem-coherent-algebraic}
Let $S$ be a scheme. Let $f : X \to B$ be a morphism of algebraic spaces
over $S$. Assume that $f$ is of finite presentation, separated, and
flat\footnote{This assumption is not necessary. See
Section \ref{section-not-flat}.}. Then $\Cohstack_{X/B}$ is
an algebraic stack over $S$.
\end{theorem}

\begin{proof}
Set $\mathcal{X} = \Cohstack_{X/B}$. We have seen that $\mathcal{X}$
is a stack in groupoids over $(\Sch/S)_{fppf}$ with diagonal representable
by algebraic spaces
(Lemmas \ref{lemma-coherent-stack} and \ref{lemma-coherent-diagonal}).
Hence it suffices to find a scheme $W$ and a surjective and smooth
morphism $W \to \mathcal{X}$.

\medskip\noindent
Let $B'$ be a scheme and let $B' \to B$ be a surjective \'etale morphism.
Set $X' = B' \times_B X$ and denote $f' : X' \to B'$ the projection.
Then $\mathcal{X}' = \Cohstack_{X'/B'}$ is equal to the $2$-fibre
product of $\mathcal{X}$ with the category fibred in sets
associated to $B'$ over the category fibred in sets associated to $B$
(Remark \ref{remark-coherent-base-change}). By the material in
Algebraic Stacks, Section \ref{algebraic-section-representable-properties}
the morphism $\mathcal{X}' \to \mathcal{X}$ is surjective and \'etale.
Hence it suffices to prove the result for $\mathcal{X}'$.
In other words, we may assume $B$ is a scheme.

\medskip\noindent
Assume $B$ is a scheme. In this case we may replace $S$ by $B$, see
Algebraic Stacks, Section \ref{algebraic-section-change-base-scheme}.
Thus we may assume $S = B$.

\medskip\noindent
Assume $S = B$. Choose an affine open covering $S = \bigcup U_i$.
Denote $\mathcal{X}_i$ the restriction of $\mathcal{X}$ to
$(\Sch/U_i)_{fppf}$. If we can find schemes $W_i$ over $U_i$ and
surjective smooth morphisms $W_i \to \mathcal{X}_i$, then we
set $W = \coprod W_i$ and we obtain a surjective smooth morphism
$W \to \mathcal{X}$. Thus we may assume $S = B$ is affine.

\medskip\noindent
Assume $S = B$ is affine, say $S = \Spec(\Lambda)$.
Write $\Lambda = \colim \Lambda_i$ as a filtered colimit with each $\Lambda_i$
of finite type over $\mathbf{Z}$. For some $i$ we can find
a morphism of algebraic spaces $X_i \to \Spec(\Lambda_i)$
which is of finite presentation, separated, and flat and whose base change
to $\Lambda$ is $X$. See
Limits of Spaces, Lemmas
\ref{spaces-limits-lemma-descend-finite-presentation},
\ref{spaces-limits-lemma-descend-separated-morphism}, and
\ref{spaces-limits-lemma-descend-flat}.
If we show that $\Cohstack_{X_i/\Spec(\Lambda_i)}$ is an
algebraic stack, then it follows by base change
(Remark \ref{remark-coherent-base-change} and
Algebraic Stacks, Section \ref{algebraic-section-change-base-scheme})
that $\mathcal{X}$ is an algebraic stack.
Thus we may assume that $\Lambda$ is a finite type $\mathbf{Z}$-algebra.

\medskip\noindent
Assume $S = B = \Spec(\Lambda)$ is affine of finite type over $\mathbf{Z}$.
In this case we will verify conditions (1), (2), (3), (4), and (5) of
Artin's Axioms, Lemma \ref{artin-lemma-diagonal-representable}
to conclude that $\mathcal{X}$ is an algebraic stack.
Note that $\Lambda$ is a G-ring, see
More on Algebra, Proposition \ref{more-algebra-proposition-ubiquity-G-ring}.
Hence all local rings of $S$ are G-rings. Thus (5) holds.
By Lemma \ref{lemma-coherent-defo-thy}
we have that $\mathcal{X}$ satisfies openness of versality, hence (4) holds.
To check (2) we have to verify axioms [-1], [0], [1], [2], and [3]
of Artin's Axioms, Section \ref{artin-section-axioms}.
We omit the verification of [-1] and axioms
[0], [1], [2], [3] correspond respectively to
Lemmas \ref{lemma-coherent-stack},
\ref{lemma-coherent-limits},
\ref{lemma-coherent-RS-star},
\ref{lemma-coherent-tangent-space}.
Condition (3) follows from Lemma \ref{lemma-coherent-existence}.
Finally, condition (1) is Lemma \ref{lemma-coherent-diagonal}.
This finishes the proof of the theorem.
\end{proof}
